%--------------------------------------------------------------------%
%
% Berkas utama templat LaTeX.
%
% author Petra Barus, Peb Ruswono Aryan
%
%--------------------------------------------------------------------%
%
% Berkas ini berisi struktur utama dokumen LaTeX yang akan dibuat.
%
%--------------------------------------------------------------------%

\documentclass[11pt, a4paper, onecolumn, oneside, final]{book}

%-------------------------------------------------------------------%
%
% Konfigurasi dokumen LaTeX untuk laporan tesis IF ITB
%
% @author Petra Novandi
%
%-------------------------------------------------------------------%
%
% Berkas asli berasal dari Steven Lolong
%
%-------------------------------------------------------------------%

% Ukuran kertas
\special{papersize=210mm,297mm}

% Setting margin
\usepackage[top=3cm,bottom=3cm,left=4cm,right=3cm]{geometry}

\usepackage{mathptmx}

\usepackage[bahasa]{babel}
\usepackage{csquotes}

% Format citation
\usepackage[backend=biber,style=apa,citestyle=authoryear]{biblatex}
\DeclareLanguageMapping{bahasa}{american-apa}
\DeclareFieldFormat{apacase}{#1}

\usepackage[utf8]{inputenc}
\usepackage{graphicx}
\usepackage{titling}
\usepackage{multirow}
\usepackage{pgfgantt}
\usepackage{float}
\usepackage{ragged2e}
\usepackage{blindtext}
\usepackage{sectsty}
\usepackage{chngcntr}
\usepackage{etoolbox}
\usepackage{amsmath}
\usepackage{hyperref}       % Package untuk link di daftar isi.
\usepackage{titlesec}       % Package Format judul
\usepackage{parskip}

% \captionsetup[longtable]{skip=0.1em}

% Line satu setengah spasi
\renewcommand{\baselinestretch}{1.5}

% Setting judul
\chapterfont{\centering \Large}
\titleformat{\chapter}[display]
  {\Large\centering\bfseries}
  {\chaptertitlename\ \thechapter}{0pt}
    {\Large\bfseries\uppercase}
  \titlespacing*{\chapter}{0pt}{-50pt}{40pt}

% Setting nomor pada subbsubsubbab
\setcounter{secnumdepth}{4}
\setcounter{tocdepth}{4}

\titleformat{\paragraph}
{\normalfont\normalsize\bfseries}{\theparagraph}{1em}{}
\titlespacing*{\paragraph}
{0pt}{3.25ex plus 1ex minus .2ex}{1.5ex plus .2ex}

\titlespacing*{\section} {0pt}{2.5ex plus 1ex minus .2ex}{1.3ex plus .2ex}
\titlespacing*{\subsection} {0pt}{2.25ex plus 1ex minus .2ex}{0.5ex plus .2ex}

\makeatletter

\makeatother

% Gambar
\DeclareGraphicsExtensions{.png,.jpg,.jpeg}

% Counter untuk figure dan table.
\counterwithin{figure}{section}
\counterwithin{table}{section}

% Nomor persamaan
\renewcommand{\theequation}{\arabic{equation}}

% Daftar Gambar / Tabel
\renewcommand\thefigure{\thechapter.\arabic{figure}}
\renewcommand\thetable{\thechapter.\arabic{table}}
%--------------------------------------------------------------------%
%
% Hypenation untuk Bahasa Indonesia
%
% @author Petra Barus
%
%--------------------------------------------------------------------%
%
% Secara otomatis LaTeX dapat langsung memenggal kata dalam dokumen,
% tapi sering kali terdapat kesalahan dalam pemenggalan kata. Untuk
% memperbaiki kesalahan pemenggalan kata tertentu, cara pemenggalan
% kata tersebut dapat ditambahkan pada dokumen ini. Pemenggalan
% dilakukan dengan menambahkan karakter '-' pada suku kata yang
% perlu dipisahkan.
%
% Contoh pemenggalan kata 'analisa' dilakukan dengan 'a-na-li-sa'
%
%--------------------------------------------------------------------%

\hyphenation {
	% A
	%
	a-na-li-sa
	a-pli-ka-si
  al-go-rit-ma
  %
  % B
	%
  ban-ding
	be-be-ra-pa
	ber-ge-rak
  %
  % C
	%
	ca-ri
  %
  % D
	%
	da-e-rah
	di-nya-ta-kan
	de-fi-ni-si
  %
  % E
	%
	e-ner-gi
	eks-klu-sif
  %
  % F
	%
	fa-si-li-tas
  %
  % G
	%
  GP-GPU
	ga-bung-an
  gam-bar
  %
  % H
	%
	ha-lang-an
  %
  % I
	%
	i-nduk
  %
  % J
	%
	ka-me-ra
	kua-li-tas
  %
  % K
	%
  ke-per-caya-an
  kus-tom-i-sasi
  %
  % L
	%
  la-ku
  %
  % M
	%
  me-la-ku-kan
  meng-guna-kan
  me-nun-juk-kan
  %
  % N
	%
  %
  % O
	%
  %
  % P
	%
  pe-ning-kat-an
  %
  % Q
	%
  %
  % R
	%
  %
  % S
	%
  %
  % T
	%
  %
  % U
	%
  u-sul
  %
  % V
	%
  %
  % W
	%
  %
  % X
	%
  %
  % Y
	%
  %
  % Z
	%
}


\makeatletter

\makeatother

\bibliography{references}

\begin{document}

    %Basic configuration
    \title{Pengembangan Pencocokan Paket Berbasis GPU pada SNORT NIDS}
    \date{}
    \author{
        Afrizal Fikri \\
        NIM 13513004
    }

    % \pagenumbering{roman}
    % \setcounter{page}{0}

    \frontmatter
    \clearpage
\pagestyle{empty}

\begin{center}
\smallskip

    \Large \bfseries \MakeUppercase{\thetitle}
    \vfill

    \Large Laporan Tugas Akhir I
    \vfill

    \large Disusun sebagai syarat kelulusan mata kuliah \\
    IF4091/Tugas Akhir I dan Seminar

    \vfill

    \large Oleh

    \Large \theauthor

    \vfill
    \begin{figure}[h]
        \centering
      	\includegraphics[width=0.15\textwidth]{resources/cover-ganesha}
    \end{figure}
    \vfill

    \large
    \uppercase{
        Program Studi Teknik Informatika \\
        Sekolah Teknik Elektro dan Informatika \\
        Institut Teknologi Bandung
    }

    Mei 2017

\end{center}

\clearpage

    \clearpage
\pagestyle{empty}

\begin{center}
\smallskip

    \Large \bfseries \MakeUppercase{\thetitle}
    \vfill

    \Large Laporan Tugas Akhir I
    \vfill

    \large Oleh

    \Large \theauthor

    \large Program Studi Teknik Informatika \\
    \normalsize \normalfont
    Sekolah Teknik Elektro dan Informatika \\
    Institut Teknologi Bandung \\

    \vfill
    \normalsize \normalfont
    Bandung, 18 Mei 2017 \\
    Mengetahui, \\~\\

    \vspace{0.6 cm}
    Pembimbing, \\

    \vspace{2 cm}
    \underline{Achmad Imam Kistijantoro, S.T., M.Sc., Ph.D.} \\
    NIP. 197308092006041001 \\

\end{center}
\clearpage

    % \chapter*{Lembar Pernyataan}

Dengan ini saya menyatakan bahwa:

\begin{enumerate}

    \item Pengerjaan dan penulisan Laporan Tugas Akhir ini dilakukan tanpa menggunakan bantuan yang tidak dibenarkan.
    \item Segala bentuk kutipan dan acuan terhadap tulisan orang lain yang digunakan di dalam penyusunan laporan tugas akhir ini telah dituliskan dengan baik dan benar.
    \item Laporan Tugas Akhir ini belum pernah diajukan pada program pendidikan di perguruan tinggi mana pun.

\end{enumerate}

Jika terbukti melanggar hal-hal di atas, saya bersedia dikenakan sanksi sesuai dengan Peraturan Akademik dan Kemahasiswaan Institut Teknologi Bandung bagian Penegakan Norma Akademik dan Kemahasiswaan khususnya Pasal 2.1 dan Pasal 2.2.
\vspace{15mm}

Bandung, 6 Desember 2016 \\~\\~\\
\theauthor


    \pagestyle{plain}

    % \input{chapters/abstract-id}
    % \input{chapters/abstract-en}
    % \input{chapters/forewords}

    \titleformat*{\section}{\centering\bfseries\Large\MakeUpperCase}
    
    \tableofcontents
    \addcontentsline{toc}{chapter}{DAFTAR ISI} 
    {%
        \let\oldnumberline\numberline%
        \renewcommand{\numberline}{\figurename~\oldnumberline}%
        \listoffigures%
    }
    \addcontentsline{toc}{chapter}{DAFTAR GAMBAR}  
    {%
        \let\oldnumberline\numberline%
        \renewcommand{\numberline}{\tablename~\oldnumberline}%
        \listoftables%
    }
    \addcontentsline{toc}{chapter}{DAFTAR TABEL}

    % \titleformat*{\section}{\centering\bfseries\large\MakeUpperCase}

    % \tableofcontents

    % \addcontentsline{toc}{chapter}{\listfigurename}
    % \listoffigures

    % \addcontentsline{toc}{chapter}{\listtablename}
    % \listoftables

    % \titleformat*{\section}{\bfseries\Large}
    % \pagenumbering{arabic}

    %----------------------------------------------------------------%
    % Konfigurasi Bab
    %----------------------------------------------------------------%
    \renewcommand{\chaptername}{BAB}
    \renewcommand{\thechapter}{\Roman{chapter}}
    %----------------------------------------------------------------%

    \titleformat*{\section}{\bfseries\large}
    \mainmatter
    %----------------------------------------------------------------%
    % Dafter Bab
    % Untuk menambahkan daftar bab, buat berkas bab misalnya `chapter-6` di direktori `chapters`, dan masukkan ke sini.
    %----------------------------------------------------------------%
    \chapter{Pendahuluan}

  Pada bab ini akan dibahas mengenai landasan dari pelaksanaan Tugas Akhir tentang sistem deteksi intrusi jaringan.

\section{Latar Belakang}

  Penjaminan keamanan informasi merupakan cara-cara untuk mencegah pengaksesan, penggunaan, penyebaran, pengubahan, penyalinan, atau perusakan data yang tidak berhak. Beberapa aspek yang penting untuk dijamin dalam sebuah informasi adalah \emph{confidentiality}, \emph{integrity} dan \emph{availability}. Ketiganya biasa disebut CIA Triad dan telah menjadi pedoman prinsip keamanan informasi selama 20 tahun lebih \parencite{perrin2008}. Pada tahun 2008, Donn Parker mengusulkan aspek tambahan penjaminan informasi yaitu \emph{possession}, \emph{authencity}, dan \emph{utility} \parencite{parker1998}. Parker berpendapat bahwa model CIA hanya berfokus pada aset saja tidak pada kontrol pengguna.

  Pengamanan dapat dilakukan pada 2 sisi sumber daya, yaitu \emph{tangible resources} dan \emph{intangible resources}. \emph{Tangible resources} yaitu aset yang berupa barang kasar (yang memiliki rupa). Contoh \emph{tangible resources} misal \emph{client}, \emph{server}, dan \emph{channel} komunikasi (komunikasi kabel dan nirkabel). \emph{Intangible resources} yaitu aset yang tidak memiliki rupa, yaitu isi dari informasi yang dikirimkan. Baik keamanan dari \emph{tangible resources} maupun \emph{intangible resources} saling terkait dan memiliki peran penting dalam menjamin aset yang lain \parencite{kizza2015}.

  Salah satu sumber daya yang penting untuk dilindungi yaitu \emph{server}. Banyak aspek security yang terdapat pada \emph{server} dan dapat dengan mudah diserang \parencite{owasp2013}. Solusi untuk pengamanan \emph{server} dari \emph{request} berbahaya adalah dengan menggunakan sistem deteksi intrusi jaringan (\emph{network intrusion detection system} abbr. NIDS). NIDS berfungsi mengenali mengenali kemungkinan serangan dari \emph{request} yang diterima pada \emph{server} atau \emph{client}. Ketika \emph{request} dianggap berbahaya, maka dapat dilakukan tindakan lanjutan untuk mencegah kerusakan lebih lanjut pada aset. NIDS yang melakukan tindakan preventif seperti ini disebut juga NIPS (\emph{network intrusion prevention system}).

  Salah satu contoh NIDS yang banyak digunakan adalah NIDS Snort. Snort adalah salah satu NIDS berbasis \emph{signature} yang mencocokan paket dengan pola serangan yang didefinisikan terlebih dahulu \parencite{snort}. Snort menggunakan PCRE (\emph{Perl Compatible Regular Expression}) \emph{rule} untuk pencocokan \emph{signature} paket. \emph{Signature} paket berasal dari komunitas penggiat keamanan dan akan diperbaharui jika diketahui ada pola serangan baru.

  Karena adanya peningkatan kecepatan \emph{traffic} internet dan banyaknya serangan yang terjadi, maka dibutuhkan NIDS yang mampu melakukan deteksi dengan lebih cepat. Salah satu \emph{bottleneck} dalam pengecekan paket adalah banyaknya \emph{rule} yang harus dicocokkan \parencite{pcre2007}. Sehingga peningkatan secara signifikan dapat dicapai salah satunya dengan meningkatkan kemampuan komputasi untuk pencocokan banyak paket secara paralel. Pemanfaatan \emph{multithreading} pada \emph{multicore processor} telah dikembangkan untuk mempercepat kinerja NIDS \parencite{multi2004}.

  Selain mengoptimalkan penggunaan CPU, alternatif yang dapat digunakan yaitu penggunaan \emph{general purpose} GPU (GPGPU) \parencite{4482891}. Dapat juga menggunakan prosesor yang mudah dikustomisasi seperti ASIC atau FPGA \parencite{fpga2008}. GPGPU banyak digunakan karena selain kustomisasi yang diperlukan lebih sedikit dari ASIC dan FPGA, perbandingan kinerja antara GPGPU dan ASIC atau FPGA tidak terlalu jauh tanpa melakukan kustomisasi lebih lanjut di tingkat \emph{low level} \parencite{gnort2008}. Maka, Tugas Akhir ini akan fokus untuk melakukan pembangunan NIPS berbasis GPU.

\section{Rumusan Masalah}

  Melakukan pencocokan banyak pola secara realtime terkait dengan beberapa masalah, yaitu algoritma yang digunakan, penyimpanan \emph{dictionary}, dan metode \emph{streaming} paket. Maka dari itu, Tugas Akhir ini akan difokuskan pada pembangunan sistem deteksi intrusi jaringan paralel dengan GPU. Dalam rangka pembangunan sistem tersebut, terdapat beberapa permasalahan yang menjadi fokus perhatian dari penelitian ini, yaitu: \\
  \begin{enumerate}
      \item Menganalisis faktor-faktor apa saja yang mampu meningkatkan kinerja pencocokan \emph{signature}
      \item Mencari rancangan-rancangan yang akan digunakan dan menentukan algoritma terbaik untuk pencocokan \emph{signature}
      \item Mengembangkan sebuah arsitektur yang memadai untuk pencocokan pola dengan GPU
      \item Membangun sistem pencegah intrusi jaringan berbasis GPU
  \end{enumerate}

\section{Tujuan}

  Tujuan dari pelaksanaan Tugas Akhir ini adalah sebagai berikut:
  \begin{enumerate}
      \item Mencari rancangan-rancangan yang akan digunakan dan menentukan algoritma terbaik untuk pencocokan \emph{signature}
      \item Mengembangkan sebuah arsitektur yang memadai untuk pencocokan pola dengan GPU
      \item Membangun sistem pencegah intrusi jaringan berbasis GPU
  \end{enumerate}

\section{Batasan Masalah}

  Merancang NIPS yang cepat dengan menggunakan GPU merupakan solusi yang cocok untuk menghadapi serangan yang semakin meningkat. Tugas Akhir ini akan difokuskan pada pembangunan sistem yang melakukan inspeksi paket secara cepat menggunakan komputasi paralel pada GPU. Dalam rangka pembangunan sistem tersebut, terdapat beberapa batasan yang digunakan dalam penelitian ini, yaitu:
  \begin{enumerate}
      \item NIPS dibuat memodifikasi NIDS Snort
      \item NIPS menggunakan platform GPGPU CUDA
      \item \emph{Payload} akan diambil dari log paket berupa berkas PCAP
      \item NIPS menggunakan \emph{rule} Talos yang disesuaikan
  \end{enumerate}

\section{Metodologi}

  Tahapan yang akan dilakukan selama pelaksanaan Tugas Akhir ini adalah sebagai berikut:
  \begin{enumerate}

      \item Analisis Permasalahan \\
      Pengerjaan Tugas Akhir ini diawali dengan analisis terhadap permasalahan yang ingin diteliti lebih lanjut. Pada tahap ini, studi literatur perlu dilakukan dalam rangka mempelajari hasil-hasil penelitian terkait yang telah dilakukan, berikut metodologi penelitian serta ketercapaian yang berhasil didapatkan. Selain itu, tahap ini juga menjadi langkah awal untuk merumuskan landasan teori, penarikan hipotesis, hingga akhirnya dicapai rancangan solusi yang akan digunakan untuk membangun sistem deteksi intrusi.

      \item Analisis Metode \\
      Pada tahap analisis metode, dilakukan eksplorasi yang bertujuan untuk mendapatkan metode-metode terbaik yang dapat digunakan sebagai komponen dalam pencocokan \emph{signature}. Analisis yang dilakukan mencakup metode \emph{streaming} dan preproses \emph{signature} paket, memori \emph{layout} untuk \emph{dictionary}, dan algoritma pencocokan \emph{signature}. Selain dilakukan eksperimen secara umum dengan mengacu pada penelitian-penelitian serupa yang pernah dilakukan sebelumnya, akan dilakukan pula eksperimen dengan beberapa modifikasi yang dianggap perlu.

      \item Implementasi Sistem \\
      Pembangunan sistem deteksi intrusi dilakukan dengan memanfaatkan hasil yang didapatkan dari tahapan sebelumnya, yakni metode-metode dan algoritma pencocokan pola.

      \item Evaluasi dan Penarikan Kesimpulan \\
      Pada tahap ini dilakukan evaluasi terhadap sistem deteksi intrusi yang telah dibangun. Selain itu, dilakukan juga penarikan kesimpulan yang didasari oleh hasil evaluasi pengujian sistem deteksi intrusi.

  \end{enumerate}

\section{Jadwal Pelaksanaan Tugas Akhir}

  Berikut adalah rencana kegiatan dan jadwal (dirinci sampai per minggu) mulai dari awal pelaksanaan Tugas Akhir I s.d. sidang Tugas Akhir beserta \emph{milestones} dan \emph{deliverables}. \\
  \begin{figure}[H]
\begin{ganttchart}[
  y unit title=0.6cm,
  y unit chart=1.5cm,
  x unit=1.2mm,
  vgrid={*{6}{draw=none},dotted},
  hgrid,
  title label anchor/.style={below=-1.6ex},
  title left shift=.05,
  title right shift=-.05,
  title height=1,
  bar label node/.style={text width=2cm, align=right, font=\scriptsize\RaggedLeft, anchor=east},
  incomplete/.style={fill=white},
  progress label text={},
  bar height=.4,
  group right shift=0,
  group top shift=0,
  group height=.3,
  group peaks height=.2,
  time slot format=isodate,
  time slot format/start date=2017-02-1
]{2017-03-01}{2017-05-31}

\ganttset{calendar week text=\currentweek}

%labels
\gantttitlecalendar{year, month=shortname, week=1} \\

%tasks
\ganttbar{Analisis Permasalahan}{2017-03-06}{2017-05-07} \\
\ganttbar{Analisis Metode}{2017-05-01}{2017-05-31} \\
\ganttbar{Implementasi Sistem}{2017-08-06}{2017-10-14} \\
\ganttbar{Evaluasi dan Penarikan Kesimpulan}{2017-10-14}{2018-01-10}

%relations
% \ganttlink{elem0}{elem1} %
% \ganttlink{elem1}{elem2} %
% \ganttlink{elem2}{elem3} %

\end{ganttchart}

\end{figure}

\begin{figure}[H]
\begin{ganttchart}[
  y unit title=0.6cm,
  y unit chart=1.5cm,
  x unit=1.2mm,
  vgrid={*{6}{draw=none},dotted},
  hgrid,
  title label anchor/.style={below=-1.6ex},
  title left shift=.05,
  title right shift=-.05,
  title height=1,
  bar label node/.style={text width=2cm, align=right, font=\scriptsize\RaggedLeft, anchor=east},
  incomplete/.style={fill=white},
  progress label text={},
  bar height=.4,
  group right shift=0,
  group top shift=0,
  group height=.3,
  group peaks height=.2,
  time slot format=isodate,
  time slot format/start date=2017-02-1
]{2017-06-01}{2017-08-31}

\ganttset{calendar week text=\currentweek}

%labels
\gantttitlecalendar{year, month=shortname, week=14} \\

%tasks
\ganttbar{Analisis Permasalahan}{2017-03-05}{2017-03-06} \\
\ganttbar{Analisis Metode}{2017-06-01}{2017-08-27} \\
\ganttbar{Implementasi Sistem}{2017-08-07}{2017-08-31} \\
\ganttbar{Evaluasi dan Penarikan Kesimpulan}{2017-10-14}{2018-01-10}

%relations
% \ganttlink{elem0}{elem1} %
% \ganttlink{elem1}{elem2} %
% \ganttlink{elem2}{elem3} %

\end{ganttchart}

\end{figure}

\begin{figure}[H]
\begin{ganttchart}[
  y unit title=0.6cm,
  y unit chart=1.5cm,
  x unit=1.2mm,
  vgrid={*{6}{draw=none},dotted},
  hgrid,
  title label anchor/.style={below=-1.6ex},
  title left shift=.05,
  title right shift=-.05,
  title height=1,
  bar label node/.style={text width=2cm, align=right, font=\scriptsize\RaggedLeft, anchor=east},
  incomplete/.style={fill=white},
  progress label text={},
  bar height=.4,
  group right shift=0,
  group top shift=0,
  group height=.3,
  group peaks height=.2,
  time slot format=isodate,
  time slot format/start date=2017-02-1
]{2017-09-01}{2017-11-30}

\ganttset{calendar week text=\currentweek}

%labels
\gantttitlecalendar{year, month=shortname, week=27} \\

%tasks
\ganttbar{Analisis Permasalahan}{2017-03-05}{2017-05-06} \\
\ganttbar{Analisis Metode}{2017-04-30}{2017-06-19} \\
\ganttbar{Implementasi Sistem}{2017-09-01}{2017-10-15} \\
\ganttbar{Evaluasi dan Penarikan Kesimpulan}{2017-10-16}{2017-11-30}

%relations
% \ganttlink{elem0}{elem1} %
% \ganttlink{elem1}{elem2} %
% \ganttlink{elem2}{elem3} %

\end{ganttchart}

\caption{\emph{Gantt chart} pelaksanaan tugas akhir}
\end{figure}


  Adapun penjelasan lebih rinci untuk jadwal pelaksanaan Tugas Akhir I hingga sidang Tugas Akhir dapat dilihat pada tabel berikut. \\
  \begin {table}[H]
\begin{center}
\caption {\emph{Deliverables} tahapan pengerjaan tugas akhir}
\begin{tabular}{|p{.5cm}|l|l|l|l|}

\hline
No & Kegiatan & \multicolumn{2}{|l|}{Waktu Pengerjaan} & \\
\hline

1 & \multirow{13}{2.5cm}{Studi literatur mengenai penelitian terkait} & \multirow{4}{*}{Maret 2017} & Minggu ke-1 & - \\
\cline{1-1} \cline{4-5}

2 & & & Minggu ke-2 & - \\
\cline{1-1} \cline{4-5}

3 & & & Minggu ke-3 & - \\
\cline{1-1} \cline{4-5}

4 & & & Minggu ke-4 & - \\
\cline{1-1} \cline{4-5}

5 & & & Minggu ke-5 & - \\
\cline{1-1} \cline{3-5}

6 & & \multirow{4}{*}{April 2017} & Minggu ke-6 & - \\
\cline{1-1} \cline{4-5}

7 & & & Minggu ke-7 & \emph{Draft} Bab I \\
\cline{1-1} \cline{4-5}

8 & & & Minggu ke-8 & - \\
\cline{1-1} \cline{4-5}

9 & & & Minggu ke-9 & \emph{Draft} Bab II \\
\cline{1-1} \cline{4-5}

10 & & \multirow{4}{*}{Mei 2017} & Minggu ke-10 & Revisi Bab II \\
\cline{1-1} \cline{3-5}

11 & & & Minggu ke-11 & \emph{Draft} Bab III \\
\cline{1-1} \cline{4-5}

12 & & & Minggu ke-12 & Revisi Bab III \\
\cline{1-1} \cline{4-5}

13 & & & Minggu ke-13 & Revisi Bab II dan III \\
\hline

\end{tabular}
\end{center}
\end{table}


    \chapter{Tinjauan Pustaka}

Bab Studi Literatur digunakan untuk mendeskripsikan kajian literatur yang terkait dengan persoalan tugas akhir. Tujuan studi literatur adalah:

\begin{enumerate}
    \item menunjukkan kepada pembaca adanya gap seperti pada rumusan masalah yang memang belum terselesaikan,
    \item memberikan pemahaman yang secukupnya kepada pembaca tentang teori atau pekerjaan terkait yang terkait langsung dengan penyelesaian persoalan, serta
    \item menyampaikan informasi apa saja yang sudah ditulis/dilaporkan oleh pihak lain (peneliti/Tugas Akhir/Tesis) tentang hasil penelitian/pekerjaan mereka yang sama atau mirip kaitannya dengan persoalan tugas akhir.
\end{enumerate}

\blindtext

\blindtext

\section{Dasar Teori}
Perujukan literatur dapat dilakukan dengan menambahkan entri baru di berkas. Tulisan ini merujuk pada \parencite{knuth2001art}

    \subsection{Subbab}

    \blindtext

    % Contoh penambahan gambar di dalam teks.
    \begin{figure}[h]
        \centering
        \includegraphics
            % Ukuran dari gambar. Lebar 0.8 dikali lebar teks.
            [width=0.8\textwidth]
            % Lokasi berkas gambar.
            {resources/chapter-2-infrastructure-diagram.png}
        % Label dari gambar.
        \caption{Contoh gambar}
    \end{figure}

    \subsubsection{Subsubbab}

    \blindtext

\section{Studi Terkait}
\blindtext

    \chapter{Analisis Permasalahan dan Solusi}

  % Pada bab ini akan dipaparkan analisis permasalahan yang ada dalam permasalahan pembangunan membangun sistem pencegahan intrusi berbasis GPU. Selain itu, akan dipaparkan juga solusi yang diajukan dalam bentuk langkah-langkah pekerjaan yang akan dilakukan.

  \section{Analisis Permasalahan}
  
    Sistem pencegah intrusi jaringan (NIPS) adalah sistem yang bertugas melakukan penyaringan paket yang dianggap berbahaya. NIPS bekerja dengan cara menangkap paket yang berjalan dalam jaringan untuk kemudian dianalisis. Hasil analisis kemudian ditindaklanjuti berdasarkan \emph{rule} yang telah didefinisikan. Paket dapat ditolak atau dibatasi. Dalam beberapa kasus, bahkan aliran paket dapat diblok.

    Dampak dari pengecekan paket adalah adanya \emph{overhead} waktu respons. Lama \emph{overhead} tergantung dari waktu pencocokan pola. Pada NIPS, hal ini dapat membuat hasil analisis menjadi tidak akurat. Maka diperlukan metode untuk melakukan analisis dengan cepat. Sebagai perbandingan, \emph{bandwidth} jaringan US Naval Postgraduate School sudah mencapai 20 Gbps dengan traffic rata-rata sebesar 200 Mbps per hari. Sedangkan maksimum paket yang dapat dianalisis oleh Snort tidak lebih besar dari 200 Mbps.

    Berbagai penelitian telah dilakukan untuk mempercepat NIDS dan NIPS. Desain paling awal yaitu menggunakan desain konkuren dengan \emph{multithreading} pada CPU \parencite{multi2004}. Desain ini mampu meningkatkan penggunaan utilitas \emph{thread} CPU secara drastis. Kemudian desain berbeda yang menggunakan GPU mulai diajukan oleh \textcite{gnort2008}. Hasil yang didapat mampu mempercepat sistem hingga 5x dibandingkan CPU dengan harga yang sepadan \parencite{smith2009}.

    % Pencocokan dapat dilakukan secara stateful, ataupun tidak. Pencocokan yang berbasis stateful akan lebih susah untuk dicek secara paralel.

    % Meski demikian, masih ada beberapa masalah terjadi dalam desain yang telah diajukan. Pencocokan paket masih menggunakan algoritma yang berbasis insi

    Berdasarkan pengukuran yang dilakukan oleh \textcite{kargus2012}, didapatkan bahwa sebagian besar beban analisa paket berada pada tahap pencarian string pada \emph{payload} paket. Pada tahapan ini, sebagian besar paket yang tidak terindikasi sebagai serangan akan diloloskan. Sehingga beban pencocokan pada tahap berikutnya, yaitu pencocokan \emph{option rule} akan berkurang drastis. Maka, fokus dari solusi yang akan diajukan yaitu implementasi desain yang akan mempercepat kinerja pencocokan string paket pada NIPS Snort. Desain yang dapat digunakan untuk memanfaatkan paralelitas dapat ditinjau dari berbagai sisi, seperti GPU \emph{multithreading}, \emph{asynchronous memory transfer}, dan struktur penyimpanan.

  \section{Analisis Solusi}
  
    % Berdasarkan permasalahan-permasalahan yang muncul dalam rangka membangun sistem pencegahan intrusi berbasis GPU, terdapat rancangan solusi yang diajukan untuk menyelesaikan permasalahan tersebut dalam bentuk langkah-langkah penelitian yang akan dilakukan. Langkah-langkah tersebut di antaranya adalah sebagai berikut.

    Eksperimen penggunaan GPGPU pada NIPS akan memodifikasi NIDPS Snort. Snort digunakan karena kode bersifat \emph{open-source} dan mudah untuk dikembangkan, sehingga pengembangan dan eksperimen tidak memerlukan biaya banyak dan dapat dieksplorasi secara mandiri. 

    Sedangkan pengembangan akan menggunakan platform GPGPU CUDA. Platform CUDA merupakan platform GPGPU yang dibuat pada GPU NVIDIA. Platform CUDA digunakan karena memiliki kinerja yang rata-rata lebih baik daripada OpenCL. Selain itu, dokumentasi dan contoh CUDA lebih banyak dan lebih mudah dipahami dibanding OpenCL.

    \subsection{Implementasi Algoritma Pencocokan \emph{Signature} pada CUDA}

      Algoritma yang akan digunakan adalah \emph{Parallel Failureless Aho-Corasick} (PFAC). Perbedaan algoritma ini dibandingkan algoritma Aho-Corasick yang biasa adalah algoritma ini tidak menggunakan \emph{failure function}. Sehingga tidak ada ketergantungan antar \emph{signature} dalam kamus. Hal ini dapat mengurangi jumlah sinkronisasi yang dilakukan tiap \emph{thread} sehingga kinerja yang didapatkan akan lebih baik. 

    \subsection{Transfer \emph{Payload} Paket}

      Pengiriman dari host ke device dan sebaliknya memiliki \emph{latency} yang cukup besar. Dalam pengujian pada \textcite{gnort2008}, didapatkan bahwa \emph{latency} dari transfer payload memiliki runtime lebih dari 80\% runtime total. Sehingga untuk menyembunyikan \emph{latency}, pengiriman akan dilakukan per batch. Ukuran batch yang digunakan akan diuji dalam ukuran 32MB, 64MB, 128MB, dan 256MB.

    \subsection{Struktur Penyimpanan \emph{Signature}}

      Penyimpanan \emph{signature} mempengaruhi besar memori yang digunakan. Struktur yang digunakan akan menggunakan trie terkompresi. Struktur digunakan untuk memaksimalkan \emph{locality} sehingga mengurangi \emph{latency}. Selain itu, dengan menurunkan konsumsi memori, besar buffer yang dapat digunakan meningkat. Diharapkan, \emph{latency} dari pengiriman juga menurun.

    \subsection{Pengujian dan Pengambilan Hasil}

      Ada dua jenis pengujian yang akan dilakukan, yaitu uji pencocokan dan \emph{stress-test}. Uji pencocokan akan dilakukan untuk menguji akurasi klasifikasi paket. Sedangkan \emph{stress-test} akan menguji ketahanan sistem dalam lingkungan yang \emph{real-time}.
    % \input{chapters/chapter-4}
    % \input{chapters/chapter-5}
    %----------------------------------------------------------------%

    % Daftar pustaka
    \renewcommand{\bibname}{Daftar Pustaka}
    \cleardoublepage
    \phantomsection
    \addcontentsline{toc}{chapter}{DAFTAR PUSTAKA}
    \printbibliography[title=Daftar Pustaka]
    % \bibliography{references}
    % \bibliographystyle{apa}

    % Index
    \appendix

    \cleardoublepage
    \phantomsection

    % \addcontentsline{toc}{part}{Lampiran}
    % \part*{Lampiran}

    % \input{chapters/appendix-1}
    % \input{chapters/appendix-2}

\end{document}
